\section{Step 9 Compare the completed leaf with its interface}\label{chacha-s09}

Call the quarter-round on separated words as well as the prepared numerical example. Inspect the four stores and the bytes between them. The worked register allocation helps explain this construction; an equivalent leaf may use different scratch registers.

\subsection{Expected observation}
Exactly the selected words change, preserved registers retain their full values and the return status is zero.

\Needspace{9\baselineskip}
\subsection{Check your understanding}
Must a free-writing solution use t0 through t3?

\begin{itemize}
\item \textbf{A} No; it must satisfy the observable interface.
\item \textbf{B} Yes; register names define the cipher.
\item \textbf{C} Only the first input may use different registers.
\end{itemize}
